% Fig. 4: from a published profile to Recovery (appendix B). Mirrors the RQ1 paper's Fig. 4: seven numbered steps
% in a zigzag (1-3 left to right, 4-6 right to left, 7 below), the final box filled, and a dashed path from the steps
% that set elements apart into a dashed box for what stays outside the denominator. Node and arrow styles follow
% figures/pipeline.tex. Box texts: outline-final.md, brief 6.10 (Fig. 4). Input inside a figure environment
% (sections/11-appendix.tex), which holds the caption and the label fig:reference.
\begin{tikzpicture}[
    font=\scriptsize,
    op/.style={draw, rounded corners=2pt, minimum height=1.05cm, align=center,
               inner sep=2pt, text width=3.4cm},
    final/.style={op, fill=black!8, text width=7.4cm},
    outside/.style={op, dashed},
    flow/.style={-{Stealth[length=1.5mm]}, semithick},
    side/.style={-{Stealth[length=1.5mm]}, semithick, dashed},
  ]
  % row 1: build the reference delta (left to right)
  \node[op] (s1) at (0,0)   {1.\ Pin the base release, its DCAT version and the published profile};
  \node[op] (s2) at (4,0)   {2.\ Compare shapes and obligation tables per (class, property), or extract tables by hand};
  \node[op] (s3) at (8,0)   {3.\ Drop inherited constraints, classify by kind, set loosened elements apart};
  % row 2: fix the denominator, then match (right to left)
  \node[op] (s4) at (8,-1.65) {4.\ Align atoms to the\\ elements they ask for};
  \node[op] (s5) at (4,-1.65) {5.\ Freeze the\\ reference deltas};
  \node[op] (s6) at (0,-1.65) {6.\ Match included changes:\\ strict, relaxed, by meaning};
  % row 3: the report (filled) and what stays outside the denominator (dashed)
  \coordinate (r3) at (0,-3.4);
  \node[final, anchor=west] (s7) at (s6.west |- r3)
        {7.\ Report Recovery:\\ aligned elements with a match / aligned elements;\\ Strict beside it};
  \node[outside] (out) at (8,-3.4)
        {Outside the denominator: elements no atom asks for, loosened elements, usage notes, unhandled constructs};
  % solid path
  \draw[flow] (s1) -- (s2);
  \draw[flow] (s2) -- (s3);
  \draw[flow] (s3) -- (s4);
  \draw[flow] (s4) -- (s5);
  \draw[flow] (s5) -- (s6);
  \draw[flow] (s6.south) -- (s6.south |- s7.north);
  % dashed path: elements set apart in step 3 and left unaligned in step 4
  \coordinate (bus) at ([xshift=3mm]s3.east);
  \draw[side] (s3.east) -- (bus) |- (out.east);
  \draw[semithick, dashed] (s4.east) -- (s4.east -| bus);
\end{tikzpicture}
